\scanpage{14}
\uline{Théorème} Soit $S$ préschéma loc. noeth, $F:(\mathrm{Sch})_{/S}^{\circ}\to\mathrm{Ens}$.

Pour que $F$ soit représentable par un $X$ \uline{localement quasi fini} sur $S$ et \uline{séparé} sur $S$, il \unsure{f|f\textsuperscript{t}} et \unsure{s|S} que $F$ satisfasse les conditions suivantes

1°) $F$ de nature locale

3°) $F$ ``commute aux $\varinjlim$ d'anneaux''

2°) $F$ compatible aux desc. fid. plts et qu. \unsure{cptes}
[Il le suffit pour $B\to\widehat B$, où $B$ est un anneau local \inserted{noeth} sur $S$] \deleted{\unsure{ess. de type fini et à ext. résiduelle finie}}

4°) $F$ commute aux $\varprojlim$ d'anneaux artiniens, pour $\widehat B=\varprojlim B_n$, avec $B$ comme dans 2°

5°) $F$ proreprésentable au dessus de \deleted{$\forall$ \unsure{les pts de} $S$}, \inserted{anneaux artiniens $A$ sur $S$} les anneaux $B$ qui \deleted{\unsure{re}}\inserted{pro}représentent $F$ en $A$ sont \uline{finis} sur \deleted{$\unsure{B}$}$A$ ]

6°) $F$ ``séparé'' (par critère valuatif\ldots)

\textcircled{7°} Pour \deleted{\unsure{tout}} $\forall A$ \deleted{comme dans 2°}, \inserted{anneau local noeth complet sur $S$}, et tout hom local $A\deleted{\unsure{\leftarrow}}\to B$, $B$ fini sur $A$, \deleted{et tt} $\unsure{\eta|y}\in F(B)$, \deleted{qui soit promodulaire}, \inserted{tel que $\unsure{\eta|y}$ soit modulaire,} [on verra que cel implique: pro-modulaire\ldots] et tout $\mathfrak p\in\operatorname{Spec}(B)$, $\unsure{\eta|y}_{\mathfrak p}\in F(B_{\mathfrak p})$ est modulaire. [Cf plus bas \unsure{des variantes|du contraire}, dans un cas particulier)
\margin{[en fait, \uuline{étale} suffit, \deleted{\unsure{et}} et cela signifie que $\unsure{\eta|y}:\operatorname{Spec}B\to F$ est étale]}

\textcircled{8°} [Pour $\forall$ \deleted{\unsure{tout}} schéma \deleted{\unsure{fini}} \inserted{noeth} irréductible $S'$ sur $S$, \deleted{pour} \inserted{et} $Z$ sur $S'$ muni de $\unsure{\eta|y}\in F(Z)=F_{S'}(Z)$, $z\in Z$ au dessus du pt gén. de $S'$, tel que $\unsure{\eta|y}$ soit modulaire

\margin{\uline{N.B.} On \unsure{fera résulter|peut vérifier} que les conditions 1° à 5°, [\unsure{et avec} 6°] $+$ 7° impliquent l'analogue de 7° sans hypo. de dim 1 ni \unsure{finitude}.}

\margin{Donner l'exemple-type ($S$ de dim 1) montrant que (8°) n'est pas surabondant, $\unsure{\simeq|=}$ si $F\Rightarrow$ non ramifié. Est-ce que 7° est surabondant? \uuline{Non}, si on \unsure{renforce aussi|remplace aussi} 4° au cas adique noeth. (pas artinien)}
\mspage{13}
